% 【深论模块】所属：第7章（由 chapters/revised/07-value.tex \input 引入）
\minitopic{探究过程为何具有不可替代的价值}《美诺》式问题可以被推广：若一个人已经拥有真答案，为什么还要追求知识、理解或更好的方法？因为认识生活不只评价某一时刻手中有什么，还评价主体能否继续到达真相、辨认反证并向他人说明路线。偶然得到的正确地图可能带人到达一次目的地，却不能在道路变化后指导选择；掌握路线关系的人能够修正。探究过程的价值表现为可重复、可迁移和可纠错，而不只是为结果增加装饰。

但“亲自探究”不是绝对价值。病人没有必要独立重做临床试验，学生也不必重新发现全部数学。证言可以高效传递知识，把个人精力留给真正需要判断的部分。关键不是所有结论都由自己生产，而是主体知道依赖谁、凭什么信任、何时需要越过摘要查原始证据。成熟的认识能动性是一种管理依赖的能力，而不是拒绝依赖。

\minitopic{好奇心需要由问题价值约束}好奇心常被列为认识德性，因为它推动主体寻找新证据。然而对任何信息都想知道，并不自动值得赞美。窥探隐私、追逐猎奇细节或不断刷新无关消息可以满足好奇，却挤占注意、伤害他人并破坏长期理解。问题选择本身具有规范结构：这项信息与什么目标相关，谁承担搜集成本，知道之后能够改进什么，是否存在不应越过的边界？

智慧由此不等于拥有最多答案。它还包括识别伪问题、把宽问题缩成可研究争点、在证据不足时暂停、在调查成本过高时选择代理，以及发现原目标值得修改。一个研究团队若精确回答了错误问题，技术成功也可能是认识失败。教育若只奖励快速解题，学生会学会在既定问题内优化，却不会检查题目是否遗漏重要对象。

\minitopic{承认无知是一项有内容的成就}“我不知道”可能只是放弃，也可能是高度校准的判断。后者至少说明未知发生在哪里：缺少数据、概念含混、竞争模型无法区分，还是主体不具备相关能力；还说明什么证据可能推进问题。能定位无知的人已经掌握问题边界，能够安排下一步探究。相反，以谦逊为名拒绝判断所有事项，会逃避应承担的认识责任。

制度也需要表达有结构的无知。公共机构不能只在报告末尾写“仍有不确定性”，而应列出不确定的来源、可能后果、正在采取的监测和触发修订的条件。这样做会暂时降低权威形象，却提高长期可信度。持续把确定语气当作专业标志，会迫使专家在证据改变后把正常修订伪装成从未犯错。

\minitopic{把价值理论转成可观察的评价}认识价值若只能在抽象层面赞美真理、理解和智慧，就难以指导教学。课程可以要求学生提交答案之外的产物：证据路径显示知识根据，最小反例检验概念掌握，反事实预测检验理解，修订说明检验对错误的响应，问题说明检验探究选择。同一作品可以在这些维度分别评分，避免用流畅表达遮蔽证据缺口。

评价还要允许失败留下价值。一个假设最终被反驳，若研究者事先明确预测、保存数据、诚实报告并由失败排除一类解释，仍产生认识收益；一个碰巧成功却删除失败记录的项目则不应得高分。奖励结构会反过来塑造探究行为。好的制度不只选出已经正确的人，还让参与者有理由暴露不确定、寻找反例和修正自己。

\begin{argumentbox}
“知道更多”不是单一尺度。评价个人或机构时，至少分别记录：真结论的数量与重要性；获得方法的可靠和透明；对关系的理解与迁移；承认并定位无知的能力；问题选择和注意分配；错误被发现后的修订。把这些维度保持可区分，才能看见价值冲突而不是制造一个虚假的总分。
\end{argumentbox}

\index{探究价值}
\index{问题价值}
\index{有结构的无知}
